\section{Background}
\label{sec:background}

The management of agent lifecycles in distributed AI systems introduces architectural challenges that are not adequately addressed by either flat agent coordination models or static supervisory hierarchies. As AI agents increasingly spawn sub-agents, delegate tasks laterally, and operate asynchronously under partial observability, the accountability chain between a human principal and an autonomous outcome becomes opaque, non-linear, and legally ambiguous. This section situates the Recursive Spawning mechanism within four intersecting lineages of research: agent lifecycle management in multi-agent systems, accountability and provenance in distributed systems, the distinction between fixed and mutable delegation chains, and hierarchical task decomposition in AI. It then positions the BX3 Framework as the structural substrate that addresses the unresolved gaps in each lineage.

% ---------------------------------------------------------------
\subsection{Agent Lifecycle Management in Multi-Agent Systems}
\label{sec:lifecycle}

Multi-agent systems (MAS) that support dynamic agent creation, migration, and retirement face a fundamental naming and identity problem: when an agent spawns a child, the child's identity must be globally resolvable, its parent must be cryptographically verifiable, and its relationship to the original human principal must be recoverable even after the spawning agent has itself terminated or been replaced. The Agent Lifecycle Protocol (ALP) defines standardized stages — initialization, activation, suspension, migration, and retirement — that govern agent state transitions across the network~\cite{alp2026,ualm2025}. However, ALP and analogous lifecycle frameworks address the \emph{sequencing} of agent states without specifying how accountability is maintained when a parent agent spawns a child at a point in the chain where no human principal is directly reachable.

A second gap in existing lifecycle frameworks is the absence of a mandatory exception propagation pathway. Most agent lifecycle specifications treat exceptions as events to be handled locally by the spawning agent, with no architectural requirement to escalate unresolved exceptions to a human-reachable ancestor. The result is that lifecycle-managed agents can enter terminal failure states — situations where the agent cannot continue but has no escalation path — without ever notifying the human principal who authorized the original directive. This failure mode is particularly acute in asynchronous multi-agent systems where the parent and child may execute in different execution contexts, on different physical hosts, or with different temporal availability.

The BX3 Framework addresses these gaps through its \purpose{} and \bounds{} layers, which together implement a mandatory exception propagation pathway that is defined at the architectural level and is not configurable by individual agent nodes. Every agent node carries an immutable parent pointer that defines the delegation chain; the \bounds{} layer enforces that any exception which exceeds the node's provisioned resolution capacity must be propagated upward to the parent, and that the parent must in turn propagate it if it cannot resolve it. This propagation invariant is not a configurable policy — it is a structural compulsion of the state machine, meaning that the exception pathway exists regardless of the agent's configuration, workload, or operational state.

% ---------------------------------------------------------------
\subsection{Accountability and Provenance in Distributed Systems}
\label{sec:accountability}

Accountability in distributed systems has traditionally been addressed through audit logging, role-based access control, and cryptographic attestation. The PROV-AGENT framework provides a unified provenance model for tracking AI agent interactions, defining a structured vocabulary for representing who did what, when, and under whose authority at each step in an agentic workflow~\cite{provcagent2025,provagent2025}. Similarly, the OpenExecution Provenance Specification defines a cryptographic three-layer ledger — analogous in structure to Certificate Transparency's append-only log architecture — for auditable autonomous-agent behavior~\cite{openexecution2025}. These provenance frameworks establish that accountability requires both a complete record of agent actions and the ability to attribute each action to a specific principal.

The limitation of existing provenance approaches is that they define the \emph{format} of accountability records without defining the \emph{routing} of unresolved exceptions. A provenance record that captures the fact that Agent A delegated Task T to Agent B, and that Agent B failed to complete Task T, does not by itself specify what should happen next — whether Agent B should retry, whether Agent A should be notified, or whether the human principal should be escalated. The routing of unresolved decisions is the gap that existing provenance standards do not address, because it requires an architectural commitment — not just a data format — that some named human principal is always reachable from any point in the delegation chain.

The BX3 Framework's \fact{} layer maintains the Forensic Ledger — an append-only, cryptographically sealed record that parallels the design of Certificate Transparency's log infrastructure~\cite{certificate-transparency} — ensuring that all state transitions, directive issuances, and exception propagations are permanently recorded with tamper-evident integrity. The \fact{} layer does not participate in decision logic; it records decisions made by other layers. This separation ensures that the ledger is authoritative even when the agent network is in a degraded state. The combination of the \fact{} layer's immutable audit trail and the \bounds{} layer's mandatory exception propagation defines a complete accountability architecture: every decision is recorded, and every unresolved decision reaches a human principal.

Recent work on cryptographic agent identity — including BAID (Binding Agent Identity) and AgenticIdentity — provides the verifiable identity substrate necessary for non-repudiation in multi-agent delegation chains~\cite{baid2025,agenticidentity2025}. These systems ensure that agent identities are not spoofable and that delegation actions can be cryptographically attributed to specific agents. The BX3 Framework complements these identity systems by providing the structural context within which identity-attributed actions are evaluated: an agent's identity is meaningful only insofar as it can be mapped to a defined scope (\bounds{}), a defined purpose alignment (\purpose{}), and a complete delegation chain back to a human anchor.

% ---------------------------------------------------------------
\subsection{Fixed Versus Mutable Delegation Chains}
\label{sec:delegation}

A delegation chain defines the path of authority by which a human principal's directive propagates through intermediate agents to terminal execution nodes. The central architectural question for any delegation chain is whether the chain is \emph{fixed} — meaning that the parent-child relationships are established at spawn time and cannot be modified after the child is created — or \emph{mutable} — meaning that parent-child relationships can be reconfigured after creation, including re-parenting a child to a different parent or inserting new intermediaries into an existing chain.

Mutable delegation chains offer flexibility: an agent can re-parent a stalled child to a more capable agent, redistribute work across a dynamically resized pool of workers, or reassign sub-tasks in response to changing conditions. However, mutability introduces a fundamental accountability problem. If a delegation chain can be modified after an action is taken, then the chain of causal responsibility for that action is retroactively altered. The provenance record for the action no longer accurately represents the delegation structure that existed at the time of the action, violating the non-repudiation requirement for accountable systems. This problem is not merely theoretical: in legal contexts where AI-generated decisions have consequences, retroactive chain-of-command modification would undermine the evidentiary value of the delegation record.

Fixed delegation chains avoid this problem by establishing an immutable parent pointer at the point of spawn. The parent pointer is recorded in the Forensic Ledger and cannot be modified. This design choice — which the BX3 Framework adopts — ensures that the delegation chain is a permanent, tamper-evident record of the actual authority structure that existed at the time of each action. The consequence is that flexibility must be achieved through other mechanisms — such as dynamic task redistribution at the sibling level rather than re-parenting — rather than by modifying the delegation chain itself.

The HDP (Human Delegation Provenance) Protocol and its IETF Internet-Draft specification define a cryptographic chain-of-custody model for human delegation that similarly prioritizes immutability over dynamic reconfiguration~\cite{hdp2026,hdp2025ietf,hdp2025arxiv,draft-helixar-hdp}. The WCA (Warrant Certificate Authorities) draft extends this model with auditable data provenance for AI-agent tool-call chains, establishing that delegated authority must be recorded as a warrant — a bounded, time-limited authorization — rather than as an open-ended permission grant~\cite{wca-draft}. The BX3 Framework's \bounds{} layer implements this warrant semantics directly: every delegation is bounded by a defined operational envelope that limits the scope of the delegated authority and specifies the conditions under which the delegation terminates or must be escalated.

The distinction between fixed and mutable chains also has implications for auditability. A fixed chain can be audited in its entirety by traversing the parent pointers from any terminal node back to the human anchor. This traversal is guaranteed to terminate because parent pointers are immutable and acyclic — a property that the BX3 state machine enforces architecturally. A mutable chain cannot be audited in this way, because the chain at the time of an action may differ from the chain as it exists at the time of the audit. The BX3 Framework's choice of fixed delegation chains is therefore not merely a design preference; it is the structural prerequisite for deterministic auditability.

% ---------------------------------------------------------------
\subsection{Hierarchical Task Decomposition in AI}
\label{sec:task-decomp}

Hierarchical task decomposition (HTD) is the process by which a complex directive is recursively broken into simpler sub-tasks, each assigned to an appropriate agent node, until all sub-tasks are terminal — meaning they can be executed without further decomposition. HTD is a well-established paradigm in classical AI planning and robotics, where it is implemented through AND/OR trees and similar formal structures~\cite{structuredagent2025}. In LLM-based agent systems, HTD has been studied extensively under frameworks including ReAcTree, AgentOrchestra, and the Recursive Delegation Engine (RDE), which formalize the decomposition criteria, delegation routing, and result aggregation mechanisms for hierarchical agent trees~\cite{reactree2025,agentorchestra2025,rde2025,rde2024,agentspawn2025}.

The core challenge in HTD for LLM-based agents is that task decomposition is inherently probabilistic — the quality of decomposition depends on the LLM's ability to correctly identify subtask boundaries, assess the capability requirements of each subtask, and route each subtask to an agent with appropriate capacity. Unlike classical planning systems where task decomposition is deterministic and fully specified, LLM-based decomposition introduces the possibility of mis-decomposition: a subtask may be incorrectly scoped, incorrectly delegated, or assigned to an agent whose \bounds{} do not cover the required operation. When this occurs, the resulting chain of spawned agents may diverge from the principal's intent in ways that are not recoverable through local correction.

Existing HTD frameworks address mis-decomposition primarily through result verification: after a subtask is executed, the parent agent verifies that the result is consistent with the expected outcome before aggregating it into the parent task's result. However, this verification is performed by the parent agent — itself an LLM-based system — and is therefore subject to the same probabilistic limitations as the original decomposition. The result is that verification is probabilistic, not deterministic, and a sufficiently novel mis-decomposition may produce a result that passes verification but is nonetheless incorrect or unsafe.

The BX3 Framework addresses this through the separation of verification (the \truth{} layer) from decomposition (the \purpose{} layer). The \truth{} layer does not verify by comparing the result to the parent agent's expectation; it verifies by comparing the result to the Forensic Ledger's record of what was delegated, what constraints were specified, and what state transitions occurred. This ledger-based verification is deterministic: it does not depend on the LLM's judgment about whether a result is reasonable. Instead, it checks whether the result falls within the envelope defined by the delegation's \bounds{} and whether the state transitions that produced it are consistent with the ledger's record of the directive. This architectural separation is what makes the BX3 Framework's HTD robust to mis-decomposition: verification is not probabilistic, and any result that falls outside the ledger-defined envelope triggers an automatic escalation through the bailout pathway.

% ---------------------------------------------------------------
\subsection{The BX3 Framework as Substrate for Immutable Parent Chains}
\label{sec:bx3-substrate}

The BX3 Framework defines a three-layer cognitive architecture — \purpose{}, \truth{}, and \bounds{} — operating over a shared Forensic Ledger maintained by the \fact{} layer. This architecture provides the structural substrate for immutable parent chains by establishing three architectural commitments that are not configurable by individual agent nodes.

First, every agent node is assigned an immutable parent pointer at spawn time. This pointer is recorded in the Forensic Ledger and defines the chain of delegation from the node back to the human anchor. The parent pointer cannot be modified after the node is created; the chain is append-only. This property ensures that the delegation chain for any agent node, at any point in time, is always recoverable and always reflects the actual authority structure that existed at the time of each action.

Second, the \bounds{} layer defines and enforces the operational envelope of each node — the set of actions the node is permitted to take, the conditions under which it may take them, and the exceptions that must be propagated upward when those conditions are exceeded. The \bounds{} layer does not merely permit or deny actions; it actively monitors operational state in real time and initiates escalation when the node's envelope is exceeded. This active constraint-satisfaction is what differentiates the BX3 Framework from role-based access control systems, which are declarative and do not actively monitor operational state.

Third, the bailout pathway — defined structurally within the \bounds{} layer — ensures that any exception which exceeds the node's provisioned resolution capacity is propagated upward through the parent chain until it reaches a node that can resolve it, or until it reaches the human anchor. This pathway is mandatory and architecture-level: no agent node can opt out of exception propagation, regardless of its configuration or operational state. The bailout pathway is the structural guarantee that the delegation chain is not merely a record of historical authority but an active escalation infrastructure.

The combination of these three commitments — immutable parent pointers, active \bounds{} enforcement, and mandatory exception propagation — defines the BX3 Framework's structural contribution to the recursive spawning problem. Existing frameworks define how agents spawn and what they do; the BX3 Framework defines the accountability infrastructure that ensures every spawn is authorized, every execution is bounded, and every unresolved exception reaches a human. The Recursive Spawning mechanism, introduced in Section~\ref{sec:recursive-spawning}, builds directly on this substrate: it leverages the immutable parent chain to ensure that spawned agents maintain verifiably correct parentage, and it leverages the bailout pathway to ensure that spawn failures or mis-spawns do not leave the system in an unresolvable state.